% \section{Related Work}
% \label{sec:related_work}

% This section reviews the state-of-the-art developments in automotive software architectures, lightweight object detection models, and the evaluation of AI inference on resource-constrained edge devices.

% \subsection{Adaptive AUTOSAR and Tooling Ecosystem}
% As vehicles transition toward Software-Defined Vehicle (SDV) architectures, ensuring safety and reliability for Advanced Driver Assistance Systems (ADAS) becomes increasingly critical. The AUTOSAR Adaptive Platform addresses these needs by providing a service-oriented architecture that supports dynamic deployment, high-performance computing, and scalable communication. According to AUTOSAR release R21-11 \cite{autosar_r21_11}, the platform enables hardware abstraction and interoperable services, forming a solid foundation for safety-critical and computationally intensive ADAS applications within modern SDV ecosystems.

% To facilitate the adoption of these standards, specialized toolchains have emerged. PopcornSAR, a global AUTOSAR solutions provider, offers a comprehensive development environment including \textit{AutoSAR.io}, a specialized authoring tool and design software for system modeling \cite{popcornsar}. Together with their Para-SDK and configuration tools, it streamlines the creation of Adaptive Applications, allowing developers to generate compliant C++ skeletons and proxies efficiently, thereby bridging the gap between theoretical specifications and practical deployment.

% \subsection{Lightweight Object Detection in Automotive}
% Efficient object detection is paramount for automotive perception systems. The YOLO (You Only Look Once) family of models has established itself as the standard for real-time inference. Recent research has focused on optimizing these architectures for automotive edge devices by integrating lightweight backbones. Lei et al. (2025) proposed the eMNv4-YOLO framework, which combines MobileNetv4 with YOLO to enhance target detection efficiency for robotic driving vehicle instrument reading \cite{lei2025emnv4}. Their work demonstrates that replacing heavy backbones with MobileNetv4 significantly reduces computational load while maintaining high accuracy in automotive scenarios.

% Practical implementations of such hybrid architectures have been further explored in the technical community. Recent developments in customizing the YOLOv11 architecture involve replacing its standard backbone with MobileNetv4 \cite{csdn_mobilenetv4}. This integration is achieved by modifying the model configuration (\texttt{.yaml}) to leverage MobileNetv4's universal inverted bottleneck blocks, resulting in a model that is significantly faster and more suitable for embedded deployment compared to the baseline YOLOv11n.

% \subsection{Traffic Light Detection \& Small Object Challenges}

% Detecting traffic lights presents unique challenges compared to general object detection due to the small scale of the targets and complex environmental illumination. Recent studies have highlighted the difficulty of maintaining high accuracy for small objects (fewer than $32 \times 32$ pixels) in real-time scenarios.

% Several algorithmic approaches have been proposed to address this. For instance, LNT-YOLO (2025) introduced a "High-Sample Matching" loss function specifically to improve the detection of small traffic signals in nighttime conditions. Similarly, YOLO-LLTS (2025) utilized a prior-guided enhancement module to tackle low-light traffic environments. However, while these methods achieve high accuracy through complex architectural additions (such as attention mechanisms or multi-branch feature interactions), they often increase computational overhead, making them difficult to deploy on cost-sensitive automotive edge nodes without specialized accelerators. Unlike these algorithmic-heavy approaches, our work focuses on a system-level optimization, leveraging the efficient MobileNetV4 backbone to achieve real-time performance within the standardized Adaptive AUTOSAR environment.

% \subsection{Sustainability and Performance on Edge Devices}
% Deploying deep learning models on edge devices like the Raspberry Pi requires a rigorous analysis of the trade-offs between performance and energy consumption. While edge devices are widely adopted for latency-critical applications in the Internet of Things (IoT), their energy constraints are often overlooked. 

% A recent study by \textit{On the Sustainability of AI Inferences in the Edge} \cite{edge_sustainability_2025} addresses this gap by characterizing the performance of neural networks on constrained devices, including the Raspberry Pi, NVIDIA Jetson Nano, and Google Coral. The study highlights that hardware and framework optimization, alongside external parameter tuning, are essential to balance F1 score, inference time, and power usage. This aligns with our approach of utilizing a quantized TFLite model on a Raspberry Pi 5, ensuring that the system meets the strict latency and energy requirements of an automotive sensor node.

% % -----------------------------------------------------------------------
% % COPY THE CONTENT BELOW INTO YOUR .bib FILE
% % -----------------------------------------------------------------------

% % \bibliographystyle{IEEEtran}
% % \bibliography{references}

% \section{Related Work}
% \label{sec:related_work}

% In this section, relevant studies are reviewed across three interconnected domains. First, the evolution of the AUTOSAR Adaptive Platform is examined as an enabler for modular, high-performance automotive software. Second, lightweight object detection methods are surveyed, emphasizing YOLO-based optimizations for resource-constrained edge devices. Finally, approaches to traffic light detection and edge sustainability are discussed, focusing on quantization and runtimes like LiteRT. While AUTOSAR provides the necessary infrastructure for service deployment, it does not inherently resolve the computational bottlenecks of deep learning. This creates a critical need for perception algorithms that are sufficiently lightweight to operate on general-purpose CPUs within the AUTOSAR environment, driving the exploration of the efficient architectures described below.

% \subsection{Adaptive AUTOSAR and Tooling Ecosystem}
% The AUTOSAR Adaptive Platform has been introduced to support service-oriented, high-performance automotive applications within Software-Defined Vehicle (SDV) architectures \cite{autosar_r21_11}\cite{autosar_ap_platform_design_r21_11}. By providing hardware abstraction, dynamic deployment, and scalable communication mechanisms, it enables the implementation of computation-intensive ADAS functions. To support practical development, specialized toolchains such as PopcornSAR offer authoring, configuration, and code generation tools that facilitate the deployment of AUTOSAR-compliant Adaptive Applications \cite{popcornsar}. The underlying layered software architecture further ensures modularity and portability across automotive platforms. Although optimizing the model architecture via lightweight backbones like MobileNetV4 significantly reduces computational cost, architectural improvements alone are often insufficient for specific safety-critical applications. Traffic light detection, in particular, demands rigorous real-time performance that purely architectural modifications may not fully satisfy on constrained hardware. To achieve true edge sustainability, these architectural optimizations must be coupled with specialized inference runtimes and quantization techniques that address the unique latency and energy constraints of the target application.

% \subsection{Lightweight Object Detection for Automotive Applications}
% The YOLO family has become famous for real-time object detection due to its end-to-end design and favorable accuracy–latency trade-offs \cite{redmon_2016_you}\cite{ali_2024_the}\cite{jegham_2025_yolo}\cite{lee_2024_comparative}\cite{ultralytics_2024_yolo11}. The practical feasibility of deploying YOLO models on resource-constrained IoT nodes has been explicitly demonstrated in various domains, confirming their suitability for edge applications \cite{9881356}.

% Recent efforts have focused on optimizing YOLO-based architectures for embedded and automotive scenarios. Similar optimization strategies for deep neural networks targeting resource-constrained embedded systems have also been reported in other application domains, such as smart agriculture, highlighting the cross-disciplinary demand for lightweight neural architectures \cite{9881382}. Rasheed and Zarkoosh analyzed optimization strategies for YOLOv11 to improve resource utilization \cite{rasheed_2025_yolov11}. In parallel, lightweight backbones have been explored to reduce computational overhead. Lei \textit{et al.} proposed eMNv4-YOLO, integrating MobileNetV4 to achieve improved efficiency in automotive-related detection tasks \cite{lei2025emnv4}. The effectiveness of MobileNetV4 as a lightweight backbone has also been demonstrated across a wide range of mobile and edge platforms \cite{qin_2024_mobilenetv4}\cite{chen_2025_mobilenetgdr}. Practical implementations further show that replacing standard YOLO backbones with MobileNetV4 in YOLOv11 significantly improves inference efficiency for embedded deployment \cite{csdn_mobilenetv4}.

% \subsection{Traffic Light Detection and Edge Inference Sustainability}
% Traffic light detection presents additional challenges due to the small size of targets and varying illumination conditions. While recent works propose algorithmic enhancements to improve detection accuracy, such approaches often increase architectural complexity and computational cost, limiting their suitability for cost-sensitive edge devices. Beyond algorithm design, the sustainability of AI inference on edge platforms has gained increasing attention. Recent studies systematically evaluate the trade-offs between accuracy, inference latency, and energy consumption on constrained devices, highlighting the importance of framework-level optimization and quantization \cite{edge_sustainability_2025}. Lightweight inference runtimes such as Google LiteRT have therefore been widely adopted to support efficient on-device execution through model conversion and optimization \cite{googleaiedge_2024_github}\cite{ultralytics_rpi5}. 

% Despite significant advances in lightweight model design, optimized inference runtimes, and AUTOSAR-based deployment frameworks, a clear research gap remains at their intersection. Existing studies typically address these aspects in isolation, focusing on standalone model performance or algorithmic improvements without considering standardized automotive middleware integration. This work bridges that gap by jointly synthesizing three complementary elements: the structural efficiency of the MobileNetV4 backbone, the runtime sustainability enabled by LiteRT-based quantization, and the standardized deployment capabilities of the Adaptive AUTOSAR Platform. The resulting holistic approach enables a cost-effective, CPU-based traffic light detection system that is both computationally sustainable and compliant with modern automotive software standards. To realize this integration, the following section presents the proposed methodology, which combines a YOLO detection head, a MobileNetV4 backbone, and an optimized inference pipeline tailored for edge-centric deployment.


\section{Related Work}
\label{sec:related_work}

Existing studies span three closely related domains: the AUTOSAR Adaptive Platform and its tooling ecosystem, lightweight object detection for automotive edge devices, and edge AI–enabled traffic light detection at the sensor level with a focus on inference sustainability. Despite notable advances in each area, their convergence into a unified and AUTOSAR-compliant edge perception framework remains underexplored.

\subsection{AUTOSAR Adaptive Platform and Tooling Ecosystem}
The AUTOSAR Adaptive Platform was introduced to support service-oriented and high-performance automotive applications within Software-Defined Vehicle (SDV) architectures \cite{autosar_r21_11}\cite{autosar_ap_platform_design_r21_11}. By providing hardware abstraction, dynamic deployment, and scalable communication mechanisms, it enables the implementation of computation-intensive ADAS functions. Toolchains such as PopcornSAR further facilitate practical development by supporting authoring, configuration, and code generation of AUTOSAR-compliant Adaptive Applications. 

Although AUTOSAR Adaptive Platform provides a standardized middleware foundation, it does not inherently mitigate the computational burden of deep learning inference. Consequently, perception models deployed on the platform must be sufficiently lightweight to operate on general-purpose CPUs, motivating the exploration of efficient detection architectures and optimized inference pipelines.

\subsection{Lightweight Object Detection for Automotive Applications}
% The YOLO family has been widely adopted for real-time object detection due to its end-to-end design and favorable accuracy-latency trade-offs \cite{redmon_2016_you}\cite{ali_2024_the}\cite{jegham_2025_yolo}\cite{lee_2024_comparative}\cite{ultralytics_2024_yolo11}. The feasibility of deploying YOLO-based detectors on resource-constrained IoT and embedded platforms has been demonstrated in multiple application domains \cite{9881356}\cite{9881382}. 

% Recent studies focus on further reducing computational cost through architectural optimization. Rasheed and Zarkoosh analyzed resource-efficient variants of YOLOv11 \cite{rasheed_2025_yolov11}, while Lei \textit{et al.} proposed eMNv4-YOLO by integrating MobileNetV4 as a lightweight backbone for automotive detection tasks \cite{lei2025emnv4}. The efficiency of MobileNetV4 has also been validated across diverse mobile and edge platforms \cite{qin_2024_mobilenetv4}\cite{chen_2025_mobilenetgdr}, and practical implementations confirm that replacing standard YOLO backbones with MobileNetV4 significantly improves inference efficiency for embedded deployment \cite{csdn_mobilenetv4}. However, most of these works focus on standalone model optimization without addressing standardized automotive middleware integration.

The YOLO family is widely adopted for real-time, resource-constrained object detection due to its favorable accuracy--latency trade-offs \cite{redmon_2016_you}\cite{ali_2024_the}\cite{11015428}\cite{10254054}. Prior iterations, notably YOLOv8, have established strong baselines for edge tasks by maintaining high precision alongside acceptable inference latency, often optimized through custom dataset fine-tuning and parameter pruning \cite{yolov8_10594419}\cite{yolov8_10374650}. Alternatively, the Single Shot MultiBox Detector (SSD) coupled with the MobileNetV3 backbone remains a standard benchmark for mobile vision. Highly favored for its minimal computing power requirements and efficient feature extraction via depthwise separable convolutions \cite{ssd_9824392}\cite{ssd_10303636}, this lightweight combination has proven highly effective for practical deployment on low-power single-board computers, such as the Raspberry Pi, serving as a reliable baseline for resource-constrained environments \cite{ssd_10303636}.

While these foundational models provide robust performance, recent studies emphasize further architectural optimizations to meet the strict latency and safety requirements of automotive edge devices. Replacing standard YOLO backbones with the advanced MobileNetV4 architecture has been shown to significantly improve inference efficiency for embedded deployment \cite{lei2025emnv4}\cite{ qin_2024_mobilenetv4}\cite{chen_2025_mobilenetgdr}. However, most existing works focus on standalone model optimization or hardware-specific compression without addressing standardized automotive middleware integration. Furthermore, they rarely provide direct comparisons against established baselines such as YOLOv8, YOLO11, and SSD-MobileNetV3 in a unified embedded environment.

\subsection{Traffic Light Detection and Edge Inference Sustainability}
Traffic light detection poses additional challenges due to small object size and varying illumination conditions. While algorithmic enhancements can improve detection accuracy, they often increase model complexity, limiting suitability for cost-sensitive edge devices. Beyond algorithm design, recent studies emphasize inference sustainability by analyzing the trade-offs among accuracy, latency, and energy consumption on constrained platforms \cite{11214219}. 

Lightweight inference runtimes such as Google LiteRT (TensorFlow Lite) have therefore gained attention for enabling efficient on-device execution through model conversion and quantization~\cite{11172322}\cite{10636604}\cite{9881382}\cite{9758676}. Despite advances in lightweight models, optimized runtimes, and AUTOSAR-based deployment frameworks, these aspects are typically studied in isolation.

In contrast, this work targets their intersection by jointly synthesizing the structural efficiency of the MobileNetV4 backbone, the runtime sustainability enabled by LiteRT-based quantization, and the standardized deployment capabilities of the AUTOSAR Adaptive Platform. The following section presents the proposed methodology, detailing the hybrid YOLO--MobileNetV4 architecture and the optimized edge-centric inference pipeline.

% \subsection{Traffic Light Detection and Edge Inference Sustainability}

% Traffic light detection introduces additional challenges such as small object sizes and varying illumination. Addressing these on cost-sensitive edge devices requires strictly balancing accuracy, latency, and energy consumption \cite{edge_sustainability_2025}. Consequently, lightweight inference runtimes like Google LiteRT have gained prominence for enabling efficient on-device execution via model quantization \cite{googleaiedge_2024_github, ultralytics_rpi5}. 

% Despite individual advances in lightweight algorithms, optimized runtimes, and AUTOSAR-based deployment frameworks, these aspects are typically studied in isolation. This work directly targets their intersection by jointly synthesizing the structural efficiency of the MobileNetV4 backbone, the runtime sustainability of LiteRT quantization, and the standardized deployment capabilities of the AUTOSAR Adaptive Platform. The following section details this hybrid YOLO--MobileNetV4 architecture and the proposed edge-centric inference pipeline.
